% from ega4-8-fr.tex
\oldpage[IV-3]{9}
Tout revient à démontrer que si $X$ est un $S_0$-préschéma quelconque,
l'application canonique
\begin{equation}
\operatorname{Hom}_{S_0}(X,S)\longrightarrow
\varprojlim_\lambda\operatorname{Hom}_{S_0}(X,S_\lambda)
\tag{8.2.3.1}\label{IV.8.2.3.1-fr}
\end{equation}
qui à tout $S_0$-morphisme $v:X\to S$ fait correspondre la famille
$(u_\lambda\circ v)$, est \emph{bijective}. Or, si $g:X\to S_0$ est le
morphisme structural et si on pose $\mathscr B=g_*(\mathcal O_X)$, qui est
une $\mathcal O_{S_0}$-Algèbre, l'application (8.2.3.1) s'identifie
canoniquement à (8.2.1.1) (II, 1.2.7), et la conclusion résulte donc de ce
qu'on a vu dans (8.2.1).

Les autres assertions de (8.2.3) sont conséquences du lemme général suivant~:

\begin{lemma}[8.2.4]
\label{IV.8.2.4-fr}
Soient $\mathbf C$ une catégorie, $T$ un objet de $\mathbf C$,
$\mathbf C_T$ la sous-catégorie des objets de $\mathbf C$ au-dessus de $T$.
Soit $(S_\lambda,u_{\lambda\mu})$ un système projectif dans $\mathbf C_T$~;
alors toute limite projective de ce système dans $\mathbf C_T$ est aussi
limite projective dans $\mathbf C$, et réciproquement.
\end{lemma}

Soit $f_\lambda:S_\lambda\to T$ le morphisme structural. Supposons que $S$
soit limite projective de $(S_\lambda,u_{\lambda\mu})$ dans $\mathbf C$, et
désignons par $u_\lambda:S\to S_\lambda$ les morphismes canoniques
correspondants. Considérons alors un système projectif de $T$-morphismes
$w_\lambda:Y\to S_\lambda$, où $Y\in\mathbf C_T$. Il existe par hypothèse un
morphisme unique (dans $\mathbf C$) $w:Y\to S$ tel que
$w_\lambda=u_\lambda\circ w$ pour tout $\lambda$. L'hypothèse que les
$u_\lambda$ sont des $T$-morphismes entraîne que les morphismes
$f_\lambda\circ u_\lambda:S\to T$ sont tous les mêmes, et ce morphisme $f$
fait donc de $S$ un $T$-objet. Si $g:Y\to T$ est le morphisme structural de
$Y$, on a alors
$f\circ w=f_\lambda\circ u_\lambda\circ w=f_\lambda\circ w_\lambda=g$ pour
tout $\lambda$, ce qui prouve que $w$ est un $T$-morphisme. Inversement,
supposons (avec les mêmes notations) que $S$ soit limite projective de
$(S_\lambda,u_{\lambda\mu})$ dans $\mathbf C_T$, et considérons maintenant un
système projectif de morphismes (de $\mathbf C$) $w_\lambda:Y\to S_\lambda$.
Les morphismes composés $f_\lambda\circ w_\lambda:Y\to T$ sont alors tous
les mêmes~: en effet, pour deux indices quelconques $\lambda$, $\mu$, il y a
un indice $\nu$ tel que $\lambda\leq\nu$ et $\mu\leq\nu$, d'où
$f_\nu=f_\lambda\circ u_{\lambda\nu}=f_\mu\circ u_{\mu\nu}$~; comme
$w_\lambda=u_{\lambda\nu}\circ w_\nu$ et
$w_\mu=u_{\mu\nu}\circ w_\nu$, on a
$f_\lambda\circ w_\lambda=f_\lambda\circ u_{\lambda\nu}\circ w_\nu
=f_\nu\circ w_\nu$, et on voit de même que
$f_\mu\circ w_\mu=f_\nu\circ w_\nu$. Si $g:Y\to T$ est l'unique morphisme
ainsi défini, $g$ fait de $Y$ un $T$-objet, et les $w_\lambda$ sont alors des
$T$-morphismes~; ils ont par suite une limite projective $w:Y\to S$ qui est
un $T$-morphisme, et \emph{a fortiori} un morphisme de $\mathbf C$~;
d'ailleurs la première partie du raisonnement montre que toute limite
projective $w'$ (dans $\mathbf C$) du système projectif $(w_\lambda)$ est
nécessairement aussi un $T$-morphisme, donc égale à $w$, ce qui achève de
prouver le lemme.

\begin{proposition}[8.2.5]
\label{IV.8.2.5-fr}
Sous les conditions de (8.2.2), soient $\alpha$ un élément de $L$,
$X_\alpha$ un $S_\alpha$-préschéma. Pour tout $\lambda\geq\alpha$, posons
$X_\lambda=X_\alpha\times_{S_\alpha}S_\lambda$, et pour
$\alpha\leq\lambda\leq\mu$, posons
$v_{\lambda\mu}=1_{X_\alpha}\times u_{\lambda\mu}$, de sorte que
$(X_\lambda,v_{\lambda\mu})$ est un système projectif de
$X_\alpha$-préschémas, dont l'ensemble d'indices est formé des
$\lambda\geq\alpha$ dans $L$. Posons de même
$X=X_\alpha\times_{S_\alpha}S$ et $v_\lambda=1_{X_\alpha}\times u_\lambda$.
Alors les $X_\alpha$-morphismes $v_\lambda:X\to X_\lambda$ font de $X$ une
limite projective du système projectif $(X_\lambda,v_{\lambda\mu})$ dans la
catégorie des $X_\alpha$-préschémas, ou dans la catégorie de tous les
préschémas.
\end{proposition}

Cela résultera encore du lemme général suivant~:

\begin{lemma}[8.2.6]
\label{IV.8.2.6-fr}
Soient $\mathbf C$ une catégorie dans laquelle les produits fibrés existent,
$q:T'\to T$ un morphisme de $\mathbf C$, $\mathbf C_T$ (resp.
$\mathbf C_{T'}$) la catégorie des objets de $\mathbf C$ au-dessus
